% from ega0-22-fr.tex
\oldpage[0\textsubscript{IV-1}]{209}

  L'intérêt de ce résultat est qu'il y a toujours de telles familles
  $(k_\alpha)$ pour lesquelles $[k:k_\alpha]$ soit fini (21.8.6). Nous ne
  démontrerons pas ce raffinement du critère de Zariski, dont nous n'aurons pas
  à faire usage. Dans (22.7), on va donner, pour les anneaux locaux complets,
  une variante (due aussi à Nagata) du critère de Zariski, qui se démontre
  essentiellement par la même méthode (avec des difficultés techniques un peu
  plus grandes).
\end{enumerate}
\end{remark}

\subsection{Le critère jacobien de Nagata}
\label{subsection:0.22.7-fr}

\begin{env}[22.7.1]
\label{0.22.7.1-fr}
Le critère jacobien de Nagata est l'analogue du critère jacobien de Zariski,
mais pour des anneaux quotients d'anneaux de \emph{séries formelles} sur un
corps. Nous en donnerons, comme Nagata [31], deux versions, présentées ici
comme critères de lissité formelle.
\end{env}

\begin{proposition}[22.7.2]
\label{0.22.7.2-fr}
Soient $k$ un corps, $A=k[[T_1,\ldots,T_r]]$ un anneau de séries formelles sur
$k$, muni de sa structure usuelle de $k$-algèbre, $\mathfrak q$ un idéal de
$A$, $B=A/\mathfrak q$. Soit $\mathfrak p$ un idéal premier de $A$ contenant
$\mathfrak q$. On suppose qu'il existe une sous-$k$-algèbre $C'$ de
$C=A/\mathfrak p$, isomorphe à une algèbre de séries formelles
$k[[X_1,\ldots,X_s]]$, telle que $C$ soit finie sur $C'$ et que le corps des
fractions $L$ de $C$ soit une extension séparable de
$L'=k((X_1,\ldots,X_s))$ (hypothèse toujours satisfaite si $k$ est de
caractéristique 0). Alors les conditions suivantes sont équivalentes~:
\begin{enumerate}
  \item[a)] $B_{\mathfrak p}=A_{\mathfrak p}/\mathfrak qA_{\mathfrak p}$ est
  une $k$-algèbre formellement lisse pour la topologie
  $\mathfrak p$-préadique.
  \item[b)] Il existe des $k$-dérivations $D_i$ $(1\leq i\leq m)$ de $A$ dans
  lui-même, et des éléments $f_i$ $(1\leq i\leq m)$ de $\mathfrak q$, tels
  que les images des $f_i$ dans $\mathfrak qA_{\mathfrak p}$ engendrent cet
  idéal de $A_{\mathfrak p}$, et que l'on ait
  $\det(D_i(f_j))\notin\mathfrak p$.
  \item[c)] $B_{\mathfrak p}$ est un anneau régulier.
\end{enumerate}
\end{proposition}

\begin{proof}
Le corps résiduel de $B_{\mathfrak p}$ est $k(\mathfrak p)$~; comme
$k((X_1,\ldots,X_s))$ est séparable sur $k$ (21.9.6.4), $k(\mathfrak p)$ est
séparable sur $k$ par hypothèse, et l'on sait déjà que dans ces conditions les
propriétés a) et c) sont équivalentes (19.6.4). D'ailleurs
$A_{\mathfrak p}$ est un anneau régulier ((17.1.4) et (17.3.2)), et comme son
corps résiduel $k(\mathfrak p)$ est séparable sur $k$, $A_{\mathfrak p}$ est
formellement lisse sur $k$ pour la topologie $\mathfrak p$-préadique
(19.6.4). Il résulte donc de (22.6.2, (ii)) que la condition a) équivaut au
fait que l'homomorphisme de $L$-espaces vectoriels
\begin{equation}
\label{0.22.7.2.1-fr}
  j_0:(\mathfrak q/\mathfrak q^2)\otimes_AL
  \longrightarrow\Omega^1_{A_{\mathfrak p}/k}\otimes_{A_{\mathfrak p}}L
  =\widehat\Omega^1_{A/k}\otimes_AL
  \tag{22.7.2.1}
\end{equation}
soit injectif.

Considérons d'autre part l'homomorphisme composé
\begin{equation}
\label{0.22.7.2.2-fr}
  j:(\mathfrak q/\mathfrak q^2)\otimes_AL
  \xrightarrow{j_0}\Omega^1_{A/k}\otimes_AL
  \longrightarrow\widehat\Omega^1_{A/k}\otimes_AL
  \tag{22.7.2.2}
\end{equation}
et montrons que la condition b) équivaut à dire que $j$ est injectif. Notons en
effet que
\[
  (\mathfrak q/\mathfrak q^2)\otimes_AL
  =(\mathfrak qA_{\mathfrak p}/\mathfrak q^2A_{\mathfrak p})
  \otimes_{A_{\mathfrak p}}L~;
\]
la condition b) signifie que si $F_i=j(f'_i\otimes1)$, où $f'_i$ est l'image
de $f_i$ dans
$\mathfrak qA_{\mathfrak p}/\mathfrak q^2A_{\mathfrak p}$, la matrice
$(\langle F_i,D_j\otimes1\rangle)$ est inversible~; donc cela entraîne que les
$F_i$ sont linéairement indépendants, et \emph{a fortiori} qu'il en est de
même des $f'_i\otimes1$~; mais comme ces derniers engendrent
$(\mathfrak qA_{\mathfrak p}/\mathfrak q^2A_{\mathfrak p})
\otimes_{A_{\mathfrak p}}L$, on en conclut que $j$ est injectif. Inversement,
supposons $j$ injectif, et prenons les $f_i$ tels que les $f'_i\otimes1$
forment

